 \section{Trabajos futuros}

A partir de los resultados obtenidos en las encuestas de evaluación de los módulos, así como de las observaciones realizadas durante el desarrollo del prototipo, se identificaron diversas oportunidades de mejora que orientan posibles líneas de trabajo futuras.

\begin{itemize}

    \item \textbf{Fortalecer la personalización y capacidad colaborativa del agente Aithena.}  
    Integrar un asistente más proactivo, capaz de ofrecer ejemplos, guías inmediatas y apoyo en la toma de decisiones durante la creación de metas, objetivos y tareas. Además, ampliar las capacidades del agente para que no solo genere planes completos desde cero, sino que también pueda trabajar a partir de estructuras iniciales proporcionadas por el estudiante (como metas, objetivos o tareas ya definidas) y ayudar a completarlas, reorganizarlas o enriquecerlas. Esto permitiría que Aithena funcione como un asistente colaborativo que complementa y refina las ideas del usuario, en lugar de limitarse únicamente a generar propuestas iniciales.


    \item \textbf{Integrar IA dentro del editor de roadmaps.}  
    Incorporar asistencia inteligente que permita complementar, reorganizar o refinar la estructura del roadmap directamente desde el editor visual. Esta funcionalidad incluiría sugerencias automáticas para ajustar la secuencia de aprendizaje, añadir tareas o subtareas, detectar redundancias y recomendar recursos complementarios.

    \item \textbf{Mejorar el sistema de recomendación de recursos educativos.}  
    Optimizar la etapa de recuperación de información dentro del flujo de Aithena, de modo que las sugerencias de recursos se ajusten mejor al formato preferido por cada estudiante, al nivel de complejidad requerido y a la naturaleza de sus metas u objetivos.


    \item \textbf{Extender la jerarquía del roadmap (subtareas).}  
    Permitir la creación de subtareas dentro de cada tarea, con el fin de representar procesos de aprendizaje más detallados y facilitar un seguimiento más granular del avance.

    \item \textbf{Implementar un temporizador Pomodoro integrado.}  
    Incorporar un sistema Pomodoro nativo permitiría sincronizar sesiones de estudio con las tareas registradas, generar métricas de dedicación y fortalecer la autorregulación del tiempo dentro de la plataforma, evitando el uso de herramientas externas.

    \item \textbf{Construir una base de conocimiento alineada al programa académico.}  
    Una futura línea de desarrollo consiste en integrar información del plan de estudios del programa de Ingeniería de Sistemas, permitiendo que Aithena genere recomendaciones y roadmaps ajustados al contenido real de cada asignatura.

    \item \textbf{Incrementar la variedad y dificultad progresiva en el módulo de Escucha Activa.}  
    Ampliar los modos de juego, incorporar niveles de dificultad adaptativos, permitir la repetición de sesiones, añadir audios en otros idiomas, integrar comandos de voz y ofrecer retroalimentación más detallada, enriqueciendo la práctica auditiva y favoreciendo un entrenamiento más completo.

    \item \textbf{Refinar los prompts de generación de contenidos y retroalimentación en el módulo de Escucha Activa.}  
    Mejorar las instrucciones utilizadas para la generación de audios y la construcción del feedback en cada modo de juego, con el fin de reducir la repetitividad, aumentar la variedad temática y asegurar que las explicaciones generadas por el sistema sean más claras, consistentes y ajustadas al desempeño del estudiante. Esto permitiría obtener contenidos auditivos más diversos y una retroalimentación más útil y contextualizada.


    \item \textbf{Mejorar la claridad visual y consistencia de la interfaz.}  
    Reforzar la coherencia visual entre módulos, mejorar la visibilidad de opciones configurables, facilitar el acceso a la IA y diseñar un tutorial inicial que ayude a comprender el flujo general de la plataforma.

    \item \textbf{Implementar mecanismos de \textit{fallback} para servicios de IA.}  
    Desarrollar una capa de respaldo que permita alternar entre diferentes modelos de lenguaje y sintetizadores de voz cuando se presenten fallos, interrupciones o limitaciones de uso en los servicios principales. Este mecanismo garantizaría la continuidad operativa de funcionalidades críticas (como la generación de contenidos, respuestas del agente o producción de audios) evitando que la plataforma dependa exclusivamente de un único proveedor y asegurando mayor estabilidad y resiliencia en el sistema.

    \item \textbf{Fortalecer el proceso de evaluación mediante estudios longitudinales y métodos mixtos.}  
    Desarrollar evaluaciones más robustas que incluyan instrumentos aplicados antes y después de la intervención (pre-test y post-test), muestras más amplias y diversas, y métricas objetivas de desempeño además de la percepción subjetiva. Asimismo, realizar estudios longitudinales permitiría observar cambios sostenidos en el tiempo y evaluar de manera más precisa el impacto real de los módulos en el desarrollo de habilidades de aprendizaje autónomo y escucha activa.

    \item \textbf{Unificar los módulos de habilidades blandas desarrollados en trabajos de grado previos.}  
    Consolidar en una única plataforma los distintos módulos de habilidades blandas desarrollados en proyectos anteriores, a fin de evitar soluciones aisladas o basadas en criterios individuales de diseño. Esta integración permitiría establecer una arquitectura coherente, estandarizar lineamientos de interacción y diseño, y facilitar la incorporación de nuevos módulos bajo una estructura común. De esta manera, se podría avanzar hacia una plataforma institucional unificada orientada al fortalecimiento integral de habilidades blandas dentro del programa académico.

\end{itemize}

